% The package xparse helps handle complex optional arguments to new commands
\usepackage{xparse}
\usepackage{ifmtarg}
\usepackage{textcase}
\usepackage{bm}
\usepackage{chemformula}

% These commands use the xparse package

% Symbols for useful constants
\NewDocumentCommand{\gasconst}{s}{%
	\IfBooleanTF{#1}%
	{\ensuremath{\num{8.314}}} {\ensuremath{R}}%
}
\NewDocumentCommand{\boltz}{s}{%
	\IfBooleanTF{#1}%
	{\ensuremath{\num{1.38e-23}}} {\ensuremath{k_B}}%
}
\NewDocumentCommand{\planck}{s}{%
	\IfBooleanTF{#1}%
	{\ensuremath{\num{6.63e-34}}} {\ensuremath{h}}%
}
\NewDocumentCommand{\avogadro}{s}{%
	\IfBooleanTF{#1}%
	{\ensuremath{\num{6.022e23}}} {\ensuremath{N_{A}}}%
}

% Declaration of non-standard units
\DeclareSIUnit\bar{bar}
\DeclareSIUnit\year{y}

\NewDocumentCommand{\RT}{}{\ensuremath{\gasconst T}}
\NewDocumentCommand{\kT}{}{\ensuremath{\boltz T}}

% Common decorations
\NewDocumentCommand{\cmp}{}{\ensuremath{*}}
\NewDocumentCommand{\cmplx}{}{\ensuremath{*}}
\NewDocumentCommand{\std}{}{\ensuremath{\circ}}
\NewDocumentCommand{\pmolar}{m}{\ensuremath{\overline{#1}}}

% Thermodynamic potentials
% Upper case is molar, lower case is volumetric

\NewDocumentCommand{\Gibbs}{s}{%
	\IfBooleanTF{#1}%
	{\ensuremath{\mathbb{G}}} {\ensuremath{\text{G}}}%
}
\NewDocumentCommand{\gibbs}{}{\ensuremath{\text{g}}}

\NewDocumentCommand{\Helmholtz}{s}{%
	\IfBooleanTF{#1}%
	{\ensuremath{\mathbb{A}}} {\ensuremath{\text{A}}}%
}
\NewDocumentCommand{\helmholtz}{}{\ensuremath{\text{a}}}

\NewDocumentCommand{\Entropy}{s}{%
	\IfBooleanTF{#1}%
	{\ensuremath{\mathbb{S}}} {\ensuremath{\text{S}}}%
}
\NewDocumentCommand{\entropy}{}{\ensuremath{\text{s}}}

\NewDocumentCommand{\Energy}{s}{%
	\IfBooleanTF{#1}%
	{\ensuremath{\mathbb{U}}} {\ensuremath{\text{U}}}%
}
\NewDocumentCommand{\energy}{}{\ensuremath{\text{u}}}

\NewDocumentCommand{\Work}{s}{%
	\IfBooleanTF{#1}%
	{\ensuremath{\mathbb{W}}} {\ensuremath{\text{W}}}%
}
\NewDocumentCommand{\Enthalpy}{s}{%
	\IfBooleanTF{#1}%
	{\ensuremath{\mathbb{H}}} {\ensuremath{\text{H}}}%
}

\NewDocumentCommand{\Vol}{s}{%
	\IfBooleanTF{#1}%
	{\ensuremath{\mathbb{V}}} {\ensuremath{V}}%
}
\NewDocumentCommand{\Volume}{s}{%
	\IfBooleanTF{#1}%
	{\ensuremath{\mathbb{V}}} {\ensuremath{V}}%
}

\NewDocumentCommand{\area}{s}{%
	\IfBooleanTF{#1}%
	{\ensuremath{\mathbb{A}}} {\ensuremath{A}}%
}

\NewDocumentCommand{\molecvol}{}{\ensuremath{\Omega}}

\NewDocumentCommand{\enthalpy}{}{\ensuremath{\text{h}}}
\NewDocumentCommand{\heat}{}{\ensuremath{q}}
\NewDocumentCommand{\Heat}{}{\ensuremath{Q}}

\NewDocumentCommand{\chempot}{}{\ensuremath{\mu}}
\NewDocumentCommand{\Press}{}{\ensuremath{P}}
\NewDocumentCommand{\ppress}{}{\ensuremath{p}}

\NewDocumentCommand{\Eact}{}{\ensuremath{E_{a}}}
\NewDocumentCommand{\acoeff}{}{\ensuremath{\gamma}}
\NewDocumentCommand{\actprod}{}{\ensuremath{Q}}
\NewDocumentCommand{\eqconst}{}{\ensuremath{K}}
\NewDocumentCommand{\fug}{}{\ensuremath{f}}
\NewDocumentCommand{\molfrac}{}{\ensuremath{X}}
\NewDocumentCommand{\moles}{}{\ensuremath{n}}
\NewDocumentCommand{\molecules}{}{\ensuremath{N}}
\NewDocumentCommand{\nsites}{}{\ensuremath{N}}
\NewDocumentCommand{\nvac}{}{\ensuremath{N_{\text{v}}}}

% Concentration
% If optional star token is given, puts brackets around optional string
% If optional star token is missing, ignores optional string and just types c
% Subscripts and superscripts are handled with the dec command below
\NewDocumentCommand{\conc}{sO{\cdots}}{%
	\IfBooleanTF{#1}%
	{\ensuremath{\left[#2\right]}} {\ensuremath{c}}%
}

% Activity
% If optional star token is given, puts curly braces around optional string
% If optional star token is missing, ignores optional string and just types a
% Subscripts and superscripts are handled with the dec command below
\NewDocumentCommand{\act}{sO{\cdots}}{%
	\IfBooleanTF{#1}%
	{\ensuremath{\set{#2}}} {\ensuremath{a}}%
}

\NewDocumentCommand{\oderivh}{O{}mm}{\ensuremath{\text{d}^{#1}#2/\text{d}{#3}^{#1}}}
\NewDocumentCommand{\oderivv}{O{}mm}{\ensuremath{\frac{\text{d}^{#1}#2}{\text{d}{#3}^{#1}}}}
\NewDocumentCommand{\pderivv}{O{}mm}{\ensuremath{\frac{\partial^{#1} #2}{\partial #3^{#1}}}}
\NewDocumentCommand{\pderivh}{O{}mm}{\ensuremath{\partial^{#1} #2/\partial #3^{#1}}}
\NewDocumentCommand{\rateconst}{}{\ensuremath{k}}
\NewDocumentCommand{\Energ}{}{\ensuremath{E}}
\NewDocumentCommand{\energ}{}{\ensuremath{\epsilon}}
\NewDocumentCommand{\DOS}{}{\ensuremath{\omega}}
\NewDocumentCommand{\exval}{m}{\ensuremath{\langle #1 \rangle}}
\NewDocumentCommand{\volfrac}{}{\ensuremath{\phi}}
\NewDocumentCommand{\mass}{}{\ensuremath{m}}
\NewDocumentCommand{\molarmass}{}{\ensuremath{M}}
\NewDocumentCommand{\flux}{}{\ensuremath{J}}
\NewDocumentCommand{\fluxp}{}{\ensuremath{\flux^{\prime}}}
\NewDocumentCommand{\diffus}{}{\ensuremath{D}}
\NewDocumentCommand{\dens}{}{\ensuremath{\rho}}

% Symbol for nucleation rate
\NewDocumentCommand{\nucrate}{}{\ensuremath{\Theta}}

% Notation for vectors
% The * token will make a unit vector
% The + token will indicate the magnitude of the vector
% Both tokens can be used together, but * must come before +

\NewDocumentCommand{\vect}{st+m}{%
	\IfBooleanTF{#1} {%
		\IfBooleanTF{#2} {%
			\ensuremath{\lvert \mathbf{\hat{#3}} \rvert}%
		}{%
			\ensuremath{\mathbf{\hat{#3}}}%
		}%
	}{%
		\IfBooleanTF{#2} {%
			\ensuremath{\lvert \mathbf{#3} \rvert}%
		}{%
			\ensuremath{\mathbf{#3}}%
		}%
	}%
}

% Shorthand partial derivatives
% The * token produces shorthand notation
%{ \ensuremath{#3_{\myrepeat{#2}{#4}}} }
% Both tokens can be used together, but * must come before +

\ExplSyntaxOn
\NewExpandableDocumentCommand{\myrepeat}{mm}
{
	\int_compare:nT { #1 > 0 }
	{
		#2 \prg_replicate:nn { #1 - 1 } { #2 }
	}
}
\ExplSyntaxOff

\NewDocumentCommand{\pderiv}{somm}{%
	\IfBooleanTF{#1} {%
		\IfNoValueTF{#2}
		{ \ensuremath{#3_{#4}} }
		{ \ensuremath{#3_{\myrepeat{#2}{#4}}} }
	}{%
		\IfNoValueTF{#2}
		{ \ensuremath{\frac{\partial #3}{\partial #4}} }
		{ \ensuremath{\frac{\partial^{#2} #3}{\partial #4^{#2}}} }
	}%
}

\NewDocumentCommand{\oderiv}{somm}{%
	\IfBooleanTF{#1} {%
		\IfNoValueTF{#2}
		{ \ensuremath{#3^{\prime}} }
		{ \ensuremath{#3^{\myrepeat{#2}{\prime}}} }
	}{%
		\IfNoValueTF{#2}
		{ \ensuremath{\frac{\text{d} #3}{\text{d}#4}} }
		{ \ensuremath{\frac{\text{d}^{#2} #3}{\text{d}#4^{#2}}} }
	}%
}

% Shorthand for divergence and gradient operators
\NewDocumentCommand{\diverg}{s}{%
	\IfBooleanTF{#1} {%
		\ensuremath{\text{div}\,}%
	}{%
		\ensuremath{\vect{\nabla} \cdot }%
	}%
}

\NewDocumentCommand{\grad}{s}{%
	\IfBooleanTF{#1} {%
		\ensuremath{\text{grad}\,}%
	}{%
		\ensuremath{\vect{\nabla} }%
	}%
}

\NewDocumentCommand{\curl}{s}{%
	\IfBooleanTF{#1} {%
		\ensuremath{\text{curl}\,}%
	}{%
		\ensuremath{\vect{\nabla} \times }%
	}%
}

\NewDocumentCommand{\mylaplac}{s}{%
	\IfBooleanTF{#1} {%
		\ensuremath{\div \grad\,}%
	}{%
		\ensuremath{\grad^2 }%
	}%
}


\NewDocumentCommand{\surfnrg}{}{\ensuremath{\gamma}}
\NewDocumentCommand{\Grate}{}{\ensuremath{G}}
\NewDocumentCommand{\freq}{}{\ensuremath{\nu}}

% Viscosity
\NewDocumentCommand{\visc}{}{\ensuremath{\mu}}
\NewDocumentCommand{\stress}{}{\ensuremath{\sigma}}
\NewDocumentCommand{\strain}{}{\ensuremath{\epsilon}}
\NewDocumentCommand{\sstress}{}{\ensuremath{\tau}}
\NewDocumentCommand{\sstrain}{}{\ensuremath{\gamma}}
\NewDocumentCommand{\sstrainrate}{}{\ensuremath{\dot{\gamma}}}
\NewDocumentCommand{\shearmod}{}{\ensuremath{G}}
\NewDocumentCommand{\bulkmod}{}{\ensuremath{K}}
\NewDocumentCommand{\freevol}{}{\ensuremath{V}}
\NewDocumentCommand{\tecoeff}{}{\ensuremath{\alpha}}

% Heat capacity
\NewDocumentCommand{\heatcap}{}{\ensuremath{C}}

% Mathematical spaces
\NewDocumentCommand{\Space}{sO{}}{%
	\IfBooleanTF{#1}{%
		\ensuremath{\mathbb{Z}^{#2}}
	}{%
		\ensuremath{\mathbb{R}^{#2}}
	}%
}

% Bold notation for gradient
\NewDocumentCommand{\bgrad}{}{\ensuremath{\pmb{\nabla}}}

% Function notation
\NewDocumentCommand{\fnc}{sO{}m}{%
	\IfBooleanTF{#1}{%
		\ensuremath{\widetilde{#3}_{#2}}
	}{%
		\ensuremath{\tilde{#3}_{#2}}
	}%
}

\theoremstyle{definition}

% Theorem-like environments
\newtheorem{example}{Example}
\newtheorem{definition}{Definition}
\newtheorem{remark}{Remark}

% Simulates filled in text for annotations I might want to keep invisible
% The * token will make the text invisible
% The + token will underline the text
% The default color will be aggiemaroon
% Both tokens can be used together, but * must come before +

\NewDocumentCommand{\fillin}{st+O{aggiemaroon}m}{%
\IfBooleanTF{#1} {%
	\IfBooleanTF{#2}%
	{\underline{\phantom{#4} \hspace{0.25in}}} {\phantom{#4} \hspace{0.25in}}
}{%
	\IfBooleanTF{#2}%
	{\underline{\textcolor{#3}{#4}}} {\textcolor{#3}{#4}}
}%
}

\NewDocumentCommand{\parspace}{}{\vspace{0.125truein}}
\NewDocumentCommand{\dparspace}{}{\vspace{0.25truein}}

% Operators
\NewDocumentCommand{\diff}{}{\ensuremath{\text{d}}}
\NewDocumentCommand{\idiff}{}{\ensuremath{\delta}}

\NewDocumentCommand{\xfact}{O{}}{%
	\IfValueTF{#1}{%
		\ensuremath{\num{#1}\times}%
	}{%
		\ensuremath{\times}%
	}%
}
\NewDocumentCommand{\dotp}{}{\ensuremath{\cdot}}
\NewDocumentCommand{\crossp}{}{\ensuremath{\times}}

% One-stop shop for different decorations for symbols
% Optional argument within < > brackets adds a superscript
% Optional argument within [ ] brackets adds a subscript
% Add a + if you want Delta, or add a * if you want differential
%
% Example 1:  \dec{\Gibbs} --> G
% Example 2:  \dec+{\Gibbs} --> Delta G
% Example 3:  \dec*{\Gibbs} --> d G
% Example 4:  \dec*[i]{\Gibbs} --> d G_i
% Example 5:  \dec*<\std>{\Gibbs} --> d G^o
% Example 5:  \dec+<\std>[i]{\Gibbs} --> Delta G^o_i

\NewDocumentCommand{\dec}{st+d<>d[]m}{%
\IfValueTF{#3}{%
	\IfValueTF{#4}{%
		\IfBooleanTF{#2}
		{\ensuremath{\Delta #5_{#4}^{#3}}} {%
			\IfBooleanTF{#1}
			{\ensuremath{\diff #5_{#4}^{#3}}} {\ensuremath{#5_{#4}^{#3}}}%
		}%
	}{%
		\IfBooleanTF{#2}
		{\ensuremath{\Delta #5^{#3}}} {%
			\IfBooleanTF{#1}
			{\ensuremath{\diff #5_{#4}^{#3}}} {\ensuremath{#5_{#4}^{#3}}}%
		}%
	}%
}{%
	\IfValueTF{#4}{%
		\IfBooleanTF{#2}
		{\ensuremath{\Delta #5_{#4}}} {%
			\IfBooleanTF{#1}
			{\ensuremath{\diff #5_{#4}}} {\ensuremath{#5_{#4}}}%
		}%
	}{%
		\IfBooleanTF{#2}
		{\ensuremath{\Delta #5}} {%
			\IfBooleanTF{#1}
			{\ensuremath{\diff #5}} {\ensuremath{#5}}%
		}%
	}%
}%
}

%----------------------------------------------------------------------------------
